% Generated draft: requires mathematical and source review.
\documentclass[11pt]{article}
\usepackage[margin=25mm]{geometry}
\usepackage{fontspec}
\setmainfont{Latin Modern Roman}
\usepackage{amsmath,amssymb,mathtools,unicode-math}
\setmathfont{Latin Modern Math}
\usepackage{enumitem,needspace,xurl,hyperref}
\setlength{\parindent}{0pt}
\setlength{\parskip}{5pt}
\setlength{\emergencystretch}{3em}
\newcommand{\N}{\mathbb{N}}
\newcommand{\Z}{\mathbb{Z}}
\newcommand{\Q}{\mathbb{Q}}
\newcommand{\R}{\mathbb{R}}
\newcommand{\C}{\mathbb{C}}
\newcommand{\F}{\mathbb{F}}
\newcommand{\A}{\mathbb{A}}
\newcommand{\PP}{\mathbb{P}}
\newcommand{\E}{\mathbb{E}}
\newcommand{\eps}{\varepsilon}
\newcommand{\dd}{\,\mathrm{d}}
\begin{document}
\section*{UnsolvedMath ID: 30004204}
\textbf{Source problem code:} catalogue entry
\subsection{Definitions and mathematical statement}
Let \(R=\operatorname{TMF}\) denote the spectrum of periodic topological modular forms. For a compact Lie group \(G\), let \(L=\operatorname{Lie}(G)\) be its real adjoint representation, with action induced by conjugation. Write \(E_G\) for the \(G\)-equivariant topological modular forms spectrum intended by the source. Representation desuspension means \(\Sigma^{-L}E_G=S^{-L}\wedge E_G\), where \(S^{-L}\) is the inverse of the representation sphere \(S^L\) in the relevant equivariant stable category.

Using the source's fixed-point notation and intended \(R\)-module structures, set
\[
M_G=E_G^G=\operatorname{TMF}^G,
\qquad
N_G=(\Sigma^{-L}E_G)^G.
\]
An \(R\)-module \(M\) is dualizable with dual \(N\) if there are \(R\)-linear evaluation and coevaluation maps
\[
\operatorname{ev}:N\otimes_R M\longrightarrow R,
\qquad
\operatorname{coev}:R\longrightarrow M\otimes_R N
\]
whose two triangle composites are the identity maps of \(M\) and \(N\), up to coherent homotopy. Here \(\otimes_R\) is the derived tensor product. In that case \(N\) is equivalent to the internal \(R\)-linear function spectrum \(F_R(M,R)\).

\textbf{Question.} For every compact Lie group \(G\), is \(M_G\) dualizable as an \(R\)-module, with dual equivalent to \(N_G\)? Equivalently, does the proposed module \(N_G\) admit evaluation and coevaluation maps satisfying the duality identities for \(M_G\)? The proposed identification of the dual is
\[
F_R(\operatorname{TMF}^G,\operatorname{TMF})
\simeq
(\Sigma^{-L}\operatorname{TMF})^G,
\]
together with the assertion that \(\operatorname{TMF}^G\) is dualizable.

\textbf{Unresolved qualification.} The supplied record does not specify the equivariant construction of \(\operatorname{TMF}\), the fixed-point functor denoted by \((-)^G\), or the construction of the indicated \(R\)-module structures. Representation suspension requires an appropriate equivariant setting; categorical fixed points and homotopy fixed points cannot silently be interchanged. The formulation above retains the source's intended objects conditionally on those choices and does not quantify over every possible equivariant refinement.
\subsection{Short English statement}
For every compact Lie group \(G\), the question asks whether its equivariant topological modular forms fixed points form a dualizable \(\operatorname{TMF}\)-module, with dual obtained by desuspending equivariant \(\operatorname{TMF}\) by the adjoint representation before taking fixed points. The equivariant model and fixed-point convention require clarification. The supplied dated curation labels the problem partially solved; that assessment is not independently verified here.
\subsection{Sources}
\begin{itemize}
\item Supplied local UnsolvedMath registry record, permanent ID \(30004204\), exact source problem code \texttt{catalogue entry}, titled ``Equivariant Topological Modular Forms Duality for Compact Lie Groups,'' in the Oberwolfach Reports Open Problems collection. The record gives proposed year 2019 and supplies no individual proposer.
\item Supplied source citation: \emph{Homotopy Theory (2020)}, Oberwolfach Reports, DOI 10.4171/owr/2019/36, pp.\ 2183--2256. Supplied source URL: \url{https://doi.org/10.4171/owr/2019/36}. The report title carries 2020, whereas the DOI and proposed-year metadata carry 2019; these details are retained as supplied.
\item Supplied literature source, ``original Oberwolfach report'': \url{https://doi.org/10.4171/OWR/2019/36}.
\item Supplied literature source, \emph{On Equivariant Topological Modular Forms}: \url{https://doi.org/10.1112/S0010437X23007509}.
\item Supplied literature source, \emph{Comparing Tempered and Equivariant Elliptic Cohomology}: \url{https://arxiv.org/abs/2311.07958}.
\item Literature assessment dated exactly \texttt{2026-08-22}: the supplied curation labels the problem \texttt{partially\_solved} and states that dualizability is proved for finite groups, compact abelian groups, and torus times finite groups, while the conjecture for all compact Lie groups remains open. This is supplied dated curation derived from the OWR record, not an independently verified literature assessment. The supplied background expressly requests independent verification of time-sensitive status and cited claims.
\item The supplied statement status is \texttt{corrected\_verified}. Its accompanying verification note states that mathematical content was preserved while notation, typography, and sentence structure were repaired for publication. This records the supplied editorial qualification and does not establish independent mathematical review of this draft.
\end{itemize}
\end{document}
